\scanpage{140}\mspage{14}
% archive sheet number: 139
Pour définir les $u_\alpha$, satisfaisant aux conditions précédentes, il suffit de les définir pour $\alpha=1$ ``et d'itérer'' :
\diagram{\begin{tikzcd}[ampersand replacement=\&,column sep=large,row sep=large,execute at end picture={\node[anchor=east] at ([xshift=-30pt]current bounding box.west) {(2.5.)};}]
\&X_1^{C_0}(f)\arrow[dd,"u_1(f)"]\arrow[dr,"p_1^{C_0}(f)"]\&\\
X\arrow[ur,"i_1^{C_0}(f)"]\arrow[dr,"i_1^{C_1}(f)"']\&\&Y\\
\&X_1^{C_1}(f)\arrow[ur,"p_1^{C_1}(f)"']\&
\end{tikzcd}}
Or $X_1^{C_0}(f)$ est défini comme une somme directe dans ${}_{X\backslash}M_{/Y}$, relatif à un ens. d'indices $\Delta(C_0,f)$ \uline{contenu} dans $\Delta(C_1,f)$ -- et les familles d'objets dont on prend la somme directe étant induites d'une famille \inserted{(non petite)} \deleted{définie sur} \deleted{\unsure{qui a pour ens.}} \inserted{dont l'ens. d'indices \unsure{serait|est}} $\Delta(\mathrm{Fl}_{\mathrm{coquar}}(M),f)$, savoir formé de \uline{tous} les diagrammes carrés
\diagram{\begin{tikzcd}[ampersand replacement=\&,column sep=large,row sep=large,execute at end picture={\node[anchor=east] at ([xshift=-25pt]current bounding box.west) {$D:$};\node[anchor=west,align=left] at ([xshift=35pt]current bounding box.east) {avec $i$ \unsure{coquarrable}\\de plus,};}]
A\arrow[r,phantom,"\unsure{-|\to}"description,"u"]\arrow[r,phantom,"\unsure{\cdot}"{pos=.72,below}]\arrow[d,"i"']\&X\arrow[d,"f"]\\
B\arrow[r,dash,"\unsure{v|\nu|u}"']\&Y
\end{tikzcd}}
la définition de $X_D,i_D,p_D$ \unsure{associés|associé} à $D$ ne dépend que de $D$, et non de l'ens $C_0$ ou $C_1$ de flèches coquarrables qui contient $D$. Ceci dit, la flèche $u_1(f)$
